\section{Assuring Experience and Trust: QoE, Security and Privacy}\label{sec:assure}

The previous sections optimise parts of the loop against targets. This section asks whether the result is
good enough for its users, and safe for them. Experience and trust share three properties. Both belong to
the whole loop: a frame arrives late because of the encoder, the scheduler, the transport pacing and the
reprojection together, and a pose leaks whether it travels to a renderer, a cache or a peer. Both are
judged at the human end, by what a user perceives and by what an adversary learns. And
both are usually inferred from proxies: network metrics stand in for experience, and attack accuracy on a
captured dataset for risk. The first half asks what users weigh and how experience and traffic
are modelled; the second asks what XR traffic reveals and what defences cost. \Cref{fig:assure-chain}
draws the measurement chain, the attack surfaces and the points where defences add delay.

\begin{figure}[t]
\centering
\begin{tikzpicture}[font=\sffamily\scriptsize,
  box/.style={draw=assure!70!black, fill=assure!8, rounded corners=2pt, align=center, text width=2.9cm,
    minimum height=10mm, inner sep=2pt, line width=0.6pt},
  sys/.style={draw=black!55, fill=black!3, rounded corners=2pt, align=center, text width=2.1cm,
    minimum height=24mm, inner sep=2pt, line width=0.6pt},
  adv/.style={draw=gapcolor!80!black, fill=gapcolor!6, rounded corners=2pt, align=center,
    minimum height=7mm, inner sep=3pt, line width=0.6pt},
  def/.style={draw=assure!70!black, fill=white, rounded corners=1.5pt, inner sep=1.5pt, line width=0.5pt},
  meas/.style={-{Stealth[length=1.8mm]}, line width=0.6pt, draw=black!75},
  assm/.style={-{Stealth[length=1.8mm]}, line width=0.6pt, draw=black!60, dashed},
  flow/.style={-{Stealth[length=1.8mm]}, line width=0.6pt, draw=assure!60!black},
  tap/.style={-{Stealth[length=1.8mm]}, line width=0.6pt, draw=gapcolor!85!black, densely dashed,
    preaction={draw=white, line width=2.2pt, -}},
  lab/.style={font=\sffamily\scriptsize, align=center, text=black}]
% --- measurement chain ---
\node[box] (n) at (1.6,0)  {Network metrics\\delay, jitter, loss, throughput};
\node[box] (a) at (6.0,0)  {Application metrics\\frame rate, stalls, MTP, quality switches};
\node[box] (p) at (10.4,0) {Perceived quality\\MOS, sickness,\\presence};
\node[box] (t) at (14.8,0) {Task outcome\\score, errors,\\co-presence};
\draw[meas] (n) -- node[lab, below=1pt] {pacing,\\ABR} (a);
\draw[meas] (a) -- node[lab, below=1pt] {stalls,\\fps} (p);
\draw[assm] (p) -- node[lab, below=1pt] {rarely\\measured} (t);
\draw[meas] ([xshift=4mm]n.north) .. controls +(0,0.8) and +(0,0.8) .. node[lab, fill=white, inner sep=1pt]
  {user studies with static impairments:\\loss outweighs delay; 75--90~ms RTT tolerated} ([xshift=-4mm]p.north);
\draw[assm] ([xshift=-6mm]n.north) .. controls +(0,1.9) and +(0,1.9) .. node[lab, fill=white, inner sep=1pt]
  {requirement tables: 5--20~ms derived from physiology, not tested with users} ([xshift=4mm]t.north);
% --- system and flows ---
\node[sys] (h) at (1.05,-2.95) {Headset\\pose, gaze,\\hand motion,\\camera};
\node[sys] (s) at (15.3,-2.95) {Edge / cloud\\renderer,\\state server,\\peers};
\draw[flow] (s.west |- 0,-2.3) -- (h.east |- 0,-2.3);
\draw[flow] (h.east |- 0,-2.95) -- (s.west |- 0,-2.95);
\draw[flow] (h.east |- 0,-3.6) -- (s.west |- 0,-3.6);
\node[lab, anchor=south west, inner sep=1pt] at (2.2,-2.3)  {frames (down, encrypted)};
\node[lab, anchor=south west, inner sep=1pt] at (2.2,-2.95) {pose, avatar motion (up)};
\node[lab, anchor=south west, inner sep=1pt] at (2.2,-3.6)  {viewport, camera (up)};
% defences on the flows
\node[def] (d1) at (7.1,-2.3)  {encryption, padding, jitter};
\node[def] (d2) at (7.1,-2.95) {DP noise (none for sync)};
\node[def] (d3) at (7.1,-3.6)  {camouflage, pose noise};
\draw[-{Stealth[length=1.8mm]}, line width=0.6pt, draw=gapcolor!85!black]
  (d1.north) .. controls +(0,0.7) and +(0.2,-0.9) .. node[lab, fill=white, inner sep=1pt, pos=0.55] {added delay, rarely priced} ([xshift=8mm]n.south);
% adversary and taps
\node[adv, text width=5.0cm] (x) at (11.6,-5.0) {On-path or co-located adversary};
\draw[tap] (9.5,-2.3)  -- (9.5,-4.65)  node[pos=0.78, right=0pt, lab, align=left] {app 92.4\%,\\activity 91\%};
\draw[tap] (11.2,-2.95) -- (11.2,-4.65) node[pos=0.70, right=0pt, lab, align=left] {keystrokes\\97.62\%};
\draw[tap] (12.75,-3.6)  -- (12.75,-4.65) node[pos=0.50, right=0pt, lab, align=left] {identity,\\scene type};
% legend and costs
\node[lab, align=left, anchor=north west, text width=8.3cm] at (0,-4.15) {%
  \tikz[baseline=-0.6ex]\draw[meas] (0,0)--(0.6,0); measured\enspace
  \tikz[baseline=-0.6ex]\draw[assm] (0,0)--(0.6,0); assumed\enspace
  \tikz[baseline=-0.6ex]\draw[tap] (0,0)--(0.6,0); attack (accuracy)\\[1pt]
  Also measured: AR user location from traffic, 94\%. Defence cost: \mbox{encrypting} 1400-byte packets misses
  a 0.35~ms edge budget; padding or jitter alone does not stop scene leakage.};
\end{tikzpicture}
\caption{Measurement chain and attack surfaces of a networked XR session. Solid arrows are links that user
or traffic studies have measured, dashed arrows links that are assumed; red taps show measured attack
accuracies, and defences sit on the flows where they add delay that has rarely been priced against the
frame deadline.}
\label{fig:assure-chain}
\end{figure}

\subsection{What users weigh}
Users of remote-rendered XR weigh loss and stalls above added delay of tens of milliseconds. On ALVR,
round-trip times (RTT) up to 90~ms were hardly noticed, while 8\% loss and a 35~Mbps cap degraded
experience \cite{W3095344018}. With CloudXR (N=30), QoE fell at RTT of 75~ms or more, at loss above 6\%,
or at 25~ms combined with 4\% loss \cite{W4402156718}. In social XR an RTT up to 77~ms had no significant
effect, while jitter did \cite{W4410087219}.

The headset separates two latencies. Rotation latency, the motion-to-photon (MTP) delay of turning the
head, is served locally by reprojection; network RTT reaches only interaction latency, the delay of
aiming, grabbing or seeing a peer react (\cref{sec:compute}). Isolated stalls of up to four frames went
unnoticed, and reprojection helped from eight frames (N=29) \cite{W4301184864}. This explains a contested
budget. Requirement papers derive 5--10~ms RTT and 20~ms interaction latency from physiology
\cite{W3035467795}, four to ten times below what users tolerated. No study has varied RTT and latency
compensation together, so the budget that edge placement should meet remains open.

The network moves perceived quality more than presence or sickness. Viewport-adaptive delay lowered
perceived quality but barely changed presence and left sickness almost untouched (20 observers)
\cite{W3037996690}. Fewer than 22\% of viewers reported sickness for AR point clouds, and two quality
levels were indistinguishable at 5~m on a HoloLens~2 (N=32) \cite{W4387831705}; content motion and
display, not the network, seem to drive sickness. Task performance is rarely measured beside
ratings. Group effects are tested almost only in pairs with
the impairment on one link \cite{W4415821547}; who is delayed in a larger group, and how the others judge
it, is unknown (\cref{sec:connect}).

These results share one study design: \emph{static impairments}. Every user study applies constant RTT,
independent loss or throttling; real 4G and 5G measurements enter only
afterwards, as inputs to models fitted under constant impairments \cite{W4405935991}. No study has
\emph{replayed measured network dynamics} to users. Wireless loss is bursty and correlated, so thresholds learned under independent loss may not transfer.

\subsection{How experience and traffic are modelled}
A QoE model maps what can be observed to what a user would report. Four families differ in what they
observe and what they were validated against (\cref{tab:assure-models}). \emph{Parametric standard models}
extend ITU-T P.1203 and P.1204.3 to tiled streaming \cite{W3200608537}, with P.919 as the test method for
short stimuli. \emph{Objective media metrics} are established for coding impairments: VMAF correlated best
on a public dataset \cite{W3115841380}, and colour PSNR reached $r=0.84$ for point clouds, depending on
content \cite{W4387831705}. Neither family sees interaction latency. \emph{Learned estimators} add it but
fit locally: a cloud VR gaming model reached a Pearson correlation of 0.93 \cite{W4400981668}. \emph{Network-side estimators} infer
quality from what an operator sees, such as encrypted 360-degree traffic from drive tests
\cite{W4384158049}, and are validated against formulas, not people. PERCEIVE predicted a parametric QoE
index within 4\%, but that index was never checked against ratings
\cite{W2809912305}.

\begin{table}[t]
\centering\small
\caption{Families of QoE models for networked XR. Each family is limited by what it observes and by the
ground truth it was validated against.}\label{tab:assure-models}
\setlength{\tabcolsep}{4pt}
\begin{tabularx}{\linewidth}{@{}>{\raggedright\arraybackslash}p{3.0cm}>{\raggedright\arraybackslash}X>{\raggedright\arraybackslash}X>{\raggedright\arraybackslash}X@{}}
\toprule
Family & Observes & Validated against & Limit\\
\midrule
Parametric standard (P.1203, P.1204.3 tiled) & Bitrate, resolution, stalls, quality switches & Ratings of short streaming stimuli & No interaction latency; 10~s stimuli\\
Objective media metrics (PSNR, VMAF, colour PSNR) & Decoded pixels or points against a reference & Ratings under coding impairments & Blind to timing; content dependent\\
Learned estimators & Network and application metrics, game state & Ratings on a few contents in one lab & Fitted and tested on the same small set\\
Network-side estimators & Encrypted traffic, radio and in-network counters & Parametric indices or activity labels & Never checked against users for interactive XR\\
Requirement tables & Human perception limits & Physiology, not networked user studies & 4--10$\times$ stricter than tolerated RTT\\
\bottomrule
\end{tabularx}
\end{table}

All four families lack \emph{temporal pooling}. P.919 uses 10~s stimuli, while presence needs
multi-minute tasks \cite{W4284879755}, and no model pools sub-second network events into the QoE of a
session.

\emph{Measurements of commercial platforms} show that the application stack, not the link, often sets what users see. WebRTC pacing in 5.56~ms batches set frame
delay at 6--21~ms over Wi-Fi~6 with an RTT of 1.5~ms \cite{W4391506055}. Under 27~Mbps throttling,
adaptive bitrate kept cloud VR frame loss at or below 5.12\% where a fixed 40~Mbps stream lost 32.51\%
\cite{W4210470894}: adaptation converts a throughput shortfall into image quality rather than lost frames.
Horizon Workrooms used 2--3~Mbps downlink and 0.6~Mbps uplink \cite{W4294643390}, two orders of magnitude
below the per-user rates in metaverse visions (\cref{sec:provision}). \emph{Statistical traffic models}
capture constant-bitrate video well, but they are open loop. The open RiftCat model matched mean but not
95th-percentile delay \cite{W3192844644}, and the 3GPP baseline underestimated effective VR bitrate by
25\% \cite{W4399340506}. None couples user motion with rate adaptation and congestion control, so
schedulers evaluated on them misjudge the tail.

\subsection{What XR traffic reveals}
XR transmits the body. Pose, gaze, viewport requests and camera frames are control signals and, at the
same time, behavioural and biometric data. Encryption hides content but not size, timing or
recipients.

\emph{Traffic metadata} reveals place and activity, an established result.
Traffic patterns of a mobile AR application located users with 94\% accuracy \cite{W2946752973}. On Meta
Quest Pro, off-the-shelf classifiers identified applications with 92.4\% accuracy and activities with
91\%, from under 10 minutes of traffic each \cite{W4406880464}. \emph{Motion as a
biometric} is the sharpest channel. Avatar motion received by an adversary's own client in multi-user VR
yielded 97.62\% of typed keystrokes across four applications \cite{W4398796475}, so state
synchronisation, which broadcasts head and hand poses to every peer, leaks by design
(\cref{sec:connect}). \emph{Viewport requests} identify users because prediction-driven tile requests
follow individual head motion (\cref{sec:anticipate}); the leakage persists under federated prediction
once prediction errors or QoE are uploaded \cite{W4223932369}. \emph{The camera uplink} of offloading AR
glasses reveals the type of physical scene through packet sizes and timing \cite{W7215028895}, so the
split chosen in \cref{sec:compute} is also a privacy decision.

\subsection{Defences and their price}
Every defence spends a resource the loop is short of: delay, bandwidth or prediction accuracy. The price
is rarely measured.

\emph{Encryption} is cheap for small packets and expensive at the tightest budgets: encrypting 1400-byte
packets did not fit a 0.35~ms edge budget, while 256-byte packets did \cite{W2923086750}.
\emph{Obfuscation} hides metadata at a cost in bandwidth or quality. Padding or jitter alone did not stop
scene leakage, and combined defences conflicted with QoS \cite{W7215028895}. Camouflaged tile requests
need more resources and degrade prediction \cite{W3207744807}. Noise remapping reduced the QoE loss of
privacy noise by 55--100\%, to 0--3\% at maximum privacy \cite{W4396941592}, but that loss is computed
from prediction accuracy; no user has been asked.

\emph{Authentication} schemes for
metaverse avatars are proven formally and costed in operation counts
\cite{W4295788720}; none is measured during an edge or platform hand-off, where re-authentication falls
on the interaction path (\cref{sec:provision}). \emph{Detection} rarely sees XR traffic: ``metaverse''
intrusion detectors report 96--99\% accuracy on IoT and enterprise datasets that contain no XR flows
\cite{W4388011607}. \emph{Blockchain} is contested: a 10-user pilot kept
88~FPS at 142 transactions per second \cite{W7161842864}, while most ledger frameworks report no latency
at all. The question is whether a ledger write sits on the interaction path.

The privacy--QoE trade-off is therefore modelled, not perceived. No testbed study reports a defence's
latency together with QoE or frame rate. A defence should be reported as a point on a privacy--QoE plane:
attacker accuracy against added 99th-percentile frame latency and bandwidth, on a real split-rendering
pipeline.

\begin{openproblems}
\item \textbf{Do network-side estimators predict what XR users report?} Estimators target parametric
  indices, and none addresses interactive remote rendering. A public benchmark pairing encrypted packet
  traces, client frame logs and per-minute ratings for cloud VR would let them be compared against people.
\item \textbf{Do thresholds from static impairments hold under real wireless dynamics?} Users have only
  seen constant RTT and independent loss. Replaying measured Wi-Fi and 5G traces, with reprojection on and
  off, would fix the latency budget for remote-rendered XR.
\item \textbf{How do sub-second events pool into session and group QoE?} No model links bursts of stalls
  to a multi-minute session, and groups beyond pairs are untested. Scripted burst schedules and asymmetric
  impairments in groups of four to six would answer both.
\item \textbf{Can a closed-loop traffic model reproduce the tail?} Open-loop models err at the 95th
  percentile. A generator coupling pose, rate adaptation and pacing, validated on several real headsets,
  would make scheduler studies credible.
\item \textbf{What does a defence cost against the frame deadline?} Measuring TLS, DTLS and QUIC with
  padding and jitter on split rendering, as attacker accuracy against 99th-percentile frame latency, would
  give providers a price; the same holds for re-authentication during hand-off.
\item \textbf{Can state synchronisation share avatars without sharing biometrics?} Broadcast poses reveal
  keystrokes. A protocol that perturbs or coarsens remote poses, tested for keystroke inference,
  co-presence and sickness, would close the channel.
\end{openproblems}
